\section{Data sources and reproducibility}\label{app:repro}

\begin{table}[!htbp]
\centering\small
\caption{Data sources. Retrieval dates and SHA-256 checksums of every file are in the
repository (\texttt{data/SOURCES.md}, \texttt{data/manifest.csv}).}
\label{tab:sources}
\begin{tabularx}{\textwidth}{@{}lXl@{}}
\toprule
Source & Use & Licence \\
\midrule
Open Food Facts, Open Beauty Facts \citep{off2026} & Product text; weak labels; gold frame & ODbL \\
TPSO open API \citep{tpso2026api} & Monthly item and aggregate indices & Thai open government data \\
TPSO monthly CPI release & Price-updated weights of the aggregates & Thai government publication \\
TPSO base-year note \citep{tpso2025cpinote} & Division weights; base-year changes & Thai government publication \\
MOC open data \citep{moc2026specs} & Official specifications (official labels) & Thai open government data \\
COICOP 2018 \citep{unsd2018coicop} & Class codes of the in-scope items & Cited \\
\bottomrule
\end{tabularx}
\end{table}

The repository rebuilds every number in this paper from public sources. \texttt{make data}
downloads the sources and checks their checksums; \texttt{make models} computes the
embeddings, trains the classifiers and runs the LLMs (cached responses are replayed, so no
API call is made when the cache is present); \texttt{make eval econ agreement} scores the
methods on gold, runs the simulator and the label-ceiling analysis; \texttt{make figures
paper} draws the figures and builds this PDF. \texttt{cpilens report} writes every number
the paper prints into \texttt{paper/numbers.tex}, and a check in continuous integration
fails if the paper's sources contain a number typed by hand.
\texttt{results/run\_manifest.json} records package versions, model revisions and the
random seed (\CfgSeed).

\section{In-scope items and crosswalk}\label{app:items}

Table~\ref{tab:items} lists the in-scope items with their COICOP 2018 classes and both
estimated weights. The crosswalk from Open Food Facts category tags to these items is in
\texttt{taxonomy/off\_categories\_to\_thcpi.csv}, with the basis of each mapping (usually the TPSO item heading it follows).

\begin{table}[!htbp]
\centering\small
\caption{In-scope items. Weight: price-updated CPI weight, \% of headline, estimated inside
official weight cells by an equal split or by NNLS. Source: TPSO; UNSD; author's
calculations.}
\label{tab:items}
\begin{tabular}{@{}llp{4.4cm}rr@{}}
\toprule
Code & Item & COICOP 2018 & Equal & NNLS \\
\midrule
\inputrows{generated/items.tex}
\bottomrule
\end{tabular}
\end{table}

\section{LLM prompt}\label{app:prompt}

The system prompt opens with the instructions below. It then lists every label by CPI group
(code, Thai and English name, and what it includes) and the official examples, each as
\texttt{text -> label}. The product is sent as the user message. The full prompt is
generated by \texttt{cpilens.classify.llm.system\_prompt}.

\VerbatimInput[fontsize=\small,breaklines,frame=single]{generated/prompt.txt}

\section{Hyperparameters}\label{app:hyper}

\begin{table}[!htbp]
\centering\small
\caption{Hyperparameters, as read from the code.}
\label{tab:hyper}
\begin{tabular}{@{}ll@{}}
\toprule
Setting & Value \\
\midrule
$k$-NN: candidate $k$ & \CfgKs \\
Linear and MLP: optimiser, learning rate, weight decay & Adam, \CfgLR, \CfgWeightDecay \\
MLP: hidden units, dropout & \CfgMLPHidden, \CfgDropout \\
Batch size (at most); early-stopping patience & \CfgBatch; \CfgPatience \\
Fine-tuning: epochs (at most), batch, learning rate & \CfgFtEpochs, \CfgFtBatch, \CfgFtLR \\
Fine-tuning: maximum input length (tokens) & \CfgFtMaxLen \\
LLM: official examples per item; maximum output tokens & \CfgFewShot; \CfgMaxTokens \\
Cascade: candidate thresholds (the last routes every product) & \CfgThresholds \\
Official-label budgets & \CfgBudgets{} and all \\
Bootstrap resamples; simulator draws; prior weight $\alpha$ & \CfgBoot; \CfgDraws; \CfgAlpha \\
Noise variants: annual dispersion (\%) & \CfgSigmaLo, \CfgSigmaHi \\
Gap between labelling passes (hours, at least) & \CfgGapHours \\
\bottomrule
\end{tabular}
\end{table}

\section{Compute and cost ledger}\label{app:costs}

Every billed API call is logged with its tokens and cost in \texttt{results/llm\_costs.csv};
Table~\ref{tab:costs} sums it. Input tokens include prompt-cache writes and reads, which are
billed at different rates.

\begin{table}[!htbp]
\centering\small
\caption{LLM calls and cost (USD) by model and stage. Smoke and probe: pipeline checks;
eval: classification of the evaluation and development sets. Source: author's API logs.}
\label{tab:costs}
\begin{tabular}{@{}llrrrr@{}}
\toprule
Model & Stage & Calls & Input tokens & Output tokens & USD \\
\midrule
\inputrows{generated/costs.tex}
\bottomrule
\end{tabular}
\end{table}
